\chapter{The exact deterministic seed and local smooth embedding}\label{comp:D}\label{D:chapter}
\renewcommand{\ModuleTag}{12}
A finite algebraic cycle first determines every query of a specified line search. Scaled copies then produce one fixed convex objective with failing initial points approaching its optimizer.


% SOURCE_BEGIN D.span.1
\begingroup
\providecommand{\R}{}
\renewcommand{\R}{\mathbb R}
\providecommand{\B}{}
\renewcommand{\B}{\mathcal B}
\providecommand{\I}{}
\renewcommand{\I}{I}
\setcounter{section}{0}
\setcounter{theorem}{0}
\setcounter{lemma}{0}
\setcounter{corollary}{0}
\setcounter{proposition}{0}
\setcounter{remark}{0}
\setcounter{equation}{0}
% SOURCE_END D.span.1


% SOURCE_BEGIN D.span.2
\section{Main result and exact implementation}

For a differentiable current point $x$, put $g=\nabla F(x)$ and $p=-Hg$, with $H\succ0$. We use the ordinary full inverse-BFGS update
\begin{equation}\label{D:bfgs}
 s=\alpha p,\quad y=\nabla F(x+s)-g,\quad t=y^Ts,\qquad
 H_+=\left(I-\frac{sy^T}{t}\right)H\left(I-\frac{ys^T}{t}\right)+\frac{ss^T}{t}.
\end{equation}
There is no resetting, damping, rescaling of the stored matrix, memory truncation, or skipped update.

The line search is Algorithm 4.6 of Lewis--Overton~\cite{LO}: initialize a lower bracket $\ell=0$, an upper bracket $h=+\infty$, and $\alpha=1$. At each trial:
\begin{enumerate}
\item If Armijo fails, set $h=\alpha$.
\item Otherwise, if weak Wolfe fails, set $\ell=\alpha$.
\item Otherwise, accept $\alpha$ and stop.
\end{enumerate}
After a rejection, try $(\ell+h)/2$ if $h<+\infty$, and $2\ell$ otherwise. The tests are
\[
 F(x+\alpha p)<F(x)+c_1\alpha g^Tp,
 \qquad \nabla F(x+\alpha p)^Tp>c_2g^Tp.
\]
All decisions below are strict; using non-strict acceptance inequalities produces the identical orbit. Every trial point, not just every accepted point, is differentiable.

\begin{theorem}\label{D:main}
There exist four convex functions $f_i\in C^2(\R^2)$, with
\[
 0\preceq\nabla^2 f_i(x)\preceq64I \quad(i=0,1,2,3;\ x\in\R^2),
\]
such that $F=\max_i f_i$ has compact sublevel sets and the unique global minimizer $x_*=0$, where $F$ is nonsmooth. For the deterministic line search specified above, with $c_1=10^{-4}$ and $c_2=1/2$, there are starting points $x_0^{(n)}\to0$, all initialized with $H_0=I$, whose full-BFGS iterates satisfy
\[
 x_k^{(n)}\longrightarrow p_n\ne0,\qquad
 F(x_k^{(n)})\downarrow F(p_n)>0,\qquad 0\notin\partial F(p_n).
\]
Every line search terminates after at most five trials. In fact, the same function, initial points, matrices, and trial decisions work for every fixed
\[
 0<c_1\le\tfrac13,\qquad c_1<c_2<1.
\]
The entire accepted and trial trajectories can be contained in an arbitrarily small neighborhood of $0$ by choosing $n$ sufficiently large.
\end{theorem}

The accepted lengths are generated by the deterministic search:
\[
 \boxed{\alpha_{3m}=1,\qquad\alpha_{3m+1}=\tfrac14,\qquad\alpha_{3m+2}=\tfrac1{16}.}
\]
Thus the construction resolves the implementation-selection issue.


\setcounter{section}{1}
\setcounter{theorem}{1}
\setcounter{lemma}{0}
\setcounter{corollary}{0}
\setcounter{proposition}{0}
\setcounter{remark}{0}
\setcounter{equation}{1}
% SOURCE_END D.span.2


% SOURCE_BEGIN D.span.3
\section{An exactly specified algebraic seed}

Let $(c,d)$ be the real pair specified below. Define
\begin{align*}
 U&=127c+16d,& V&=16c-127d-16,& R&=c-8d,\\
 K&=8c^2+8d^2+d-8,& T&=c^2-c+d^2,& Y&=(c-1)^2+d^2.
\end{align*}
The defining equations are
\begin{equation}\label{D:polys}
 \boxed{264192RK-U V(U+129)=0,\qquad16VTK-RUY=0.}
\end{equation}
Choose the unique root in the closed $\ell_\infty$ box of radius $10^{-30}$ centered at the rational pair
\begin{align*}
 c_0&=0.31848142762711863064645131250290703717912205023605,\\
 d_0&=-1.02523211967574872548135087001300414284511909318444.
\end{align*}
The decimals here are exact terminating rationals defining the box. Set
\begin{equation}\label{D:rho}
 \rho=\frac{129R}{V(U+129)}
       \simeq0.060199206107543764815373666287.
\end{equation}

\begin{lemma}\label{D:root}
The designated box contains exactly one solution of~\eqref{D:polys}. All denominators in this chapter are nonzero there, and $0.060<\rho<0.061$.
\end{lemma}
\begin{proof}
Let $P=(P_1,P_2)^T$ denote the left sides of~\eqref{D:polys}, let $\bar r=(c_0,d_0)$, and let $M=DP(\bar r)^{-1}$, an exactly specified rational matrix. On the box $X$, exact rational interval arithmetic gives
\[
 \sup_{z\in X}\|I-MDP(z)\|_\infty<10^{-26},
 \qquad\|MP(\bar r)\|_\infty<10^{-50}.
\]
Hence $z\mapsto z-MP(z)$ is a strict contraction mapping $X$ into itself, because
\[
 10^{-50}+10^{-26}\,10^{-30}<10^{-30}.
\]
The contraction theorem gives the unique root. The interval calculation also bounds the denominators away from zero and certifies~\eqref{D:rho}. These bounds are exact rational certificate inequalities. Appendix~\ref{comp:seed_certificate} specifies the polynomial residual, coefficient bounds, sound interval operations and finite query inequalities. The contraction argument explains how these bounds imply existence and uniqueness. The formal root and coefficient certificates are separately indexed in the theorem ledger.
\end{proof}

Set
\[
 g_0=(1,0)^T,\qquad g_1=(-1,-1/8)^T,\qquad g_2=(c,d)^T,
 \qquad G(z)=\max_{0\le i\le2}g_i^Tz.
\]
Start with $M_0=I$, and define
\[
 M_1=\begin{pmatrix}129/256&-1/16\\-1/16&1\end{pmatrix},
 \qquad M_2=\B(M_1,g_1,g_2,1/4),
\]
where $\B$ denotes~\eqref{D:bfgs} with the stated current and next gradients. Then
\begin{equation}\label{D:closure}
 \B(I,g_0,g_1,1)=M_1,\qquad
 \boxed{\B(M_2,g_2,g_0,1/16)=\rho I.}
\end{equation}
The three displacements are
\begin{equation}\label{D:steps}
 s_0=(-1,0)^T,\quad s_1=(127/1024,1/64)^T,\quad
 s_2=\rho(1-c,-d)^T.
\end{equation}
Define the starting point by
\begin{equation}\label{D:x0}
 z_0=\frac1{1-\rho}
 \begin{pmatrix}897/1024-\rho(1-c)\\ \rho d-1/64\end{pmatrix},
 \qquad z_{j+1}=z_j+s_j.
\end{equation}
Then
\[
 \boxed{z_3=\rho z_0.}
\]
For clarity, approximate points are
\[
\begin{array}{c|rr}
 j& (z_j)_1&(z_j)_2\\\hline
0& 0.888432624149&-0.082297397686\\
1&-0.111567375851&-0.082297397686\\
2& 0.012456061649&-0.066672397686\\
3& 0.053482938654&-0.004954238005
\end{array}
\]
The uniquely active gradient at $z_j$ is $g_j$ for $j=0,1,2$, and $g_0$ at $z_3$.

% SOURCE_END D.span.3


% SOURCE_BEGIN D.span.4
\subsection{Derivation and independent checks of the matrix closure}
For arbitrary symbolic $(c,d)$ with the denominators nonzero, put
\[
 \beta=\frac{66048RK}{UV(U+129)},\qquad
 \gamma=\frac{VTK}{RUY},\qquad
 \widetilde\rho=\frac{2064TK}{U(U+129)Y}.
\]
Direct rational identities give
\[
 \B\bigl(\B(M_1,g_1,g_2,\beta),g_2,g_0,\gamma\bigr)
 =\widetilde\rho I.
\]
One way to derive the formulas is to write the middle update as $C+\beta D$. Requiring $(C+\beta D)g_2$ to be parallel to $g_0-g_2$ is linear in $\beta$. Its solution is the formula above. If $w\perp(g_0-g_2)$, set
\[
 \widetilde\rho=\frac{w^T(C+\beta D)w}{w^Tw}
\]
and choose $\gamma$ so that the last displacement is $\widetilde\rho(g_0-g_2)$. The last update then gives $\widetilde\rho I$: its two projection factors are the orthogonal projector onto $(g_0-g_2)^\perp$.

Equations~\eqref{D:polys} are exactly $\beta=1/4$ and $\gamma=1/16$. Also
\[
 \widetilde\rho-\rho
 =\frac{2064R}{V(U+129)}\left(\gamma-\frac1{16}\right).
\]
This proves~\eqref{D:closure} and~\eqref{D:steps}. The Lean formalization checks these exact identities from the full-BFGS formulas, together with the secant identity and the equivalence of the inverse and Hessian updates. It proves the required positive curvatures and positive definiteness, including the stated algebraic-root box.


\setcounter{section}{2}
\setcounter{theorem}{1}
\setcounter{lemma}{1}
\setcounter{corollary}{0}
\setcounter{proposition}{0}
\setcounter{remark}{0}
\setcounter{equation}{6}
% SOURCE_END D.span.4


% SOURCE_BEGIN D.span.5
\section{Every deterministic oracle decision is certified}\label{comp:trial_table}
Use $c_1=10^{-4}$ and $c_2=1/2$ for the displayed table. The line searches below use $p_j=-M_jg_j$. At each rejected trial, the objective is strictly greater than the current value, not merely greater than the Armijo threshold.
\begin{center}\small
\begin{tabular}{cccrl}\toprule
Phase & Trial $\alpha$ & Active branch & $G(z_j+\alpha p_j)$ & Decision\\\midrule
0&1&1&0.121854550561&Accept\\\midrule
1&1&0&0.384526374149&Armijo rejection\\
1&1/2&0&0.136479499149&Armijo rejection\\
1&1/4&2&0.072321707900&Accept\\\midrule
2&1&0&0.668886093720&Armijo rejection\\
2&1/2&0&0.340671077685&Armijo rejection\\
2&1/4&0&0.176563569667&Armijo rejection\\
2&1/8&0&0.094509815658&Armijo rejection\\
2&1/16&0&0.053482938654&Accept\\\bottomrule
\end{tabular}
\end{center}
The starting values in the three phases are approximately
\[
 (0.888432624149,\;0.121854550561,\;0.072321707900).
\]
The ratios at the accepted steps are
\[
\begin{array}{c|rr}
 j & \dfrac{G(z_j)-G(z_{j+1})}{-g_j^Ts_j}
 &\dfrac{g_{j+1}^Ts_j}{g_j^Ts_j}\\\hline
0&0.766578073588&-1\\
1&0.393190937094&-0.186383158092\\
2&0.375205963573&-0.817119670485
\end{array}
\]
The Lean certificates prove the stated rational enclosures and strict comparisons for the algebraic root. In particular, each accepted Armijo ratio exceeds $1/3$, and every accepted Wolfe ratio is negative. The rejected trial values exceed their respective current values. Therefore the identical trial tree works for \emph{all} $0<c_1\le1/3$ and $c_1<c_2<1$.

All nine trial points have a unique active branch. They satisfy the additional uniform bounds
\begin{equation}\label{D:trialbounds}
 |g_i^Tz_{\rm trial}|<3\quad(i=0,1,2),
 \qquad (z_{\rm trial})_2<2,
 \qquad G(z_{\rm trial})>0.
\end{equation}
The accepted points satisfy these bounds as well.

Positive homogeneity and~\eqref{D:closure}--\eqref{D:x0} prove by induction that the deterministic full-BFGS run on $G$ is
\[
 x_{3m+j}=\rho^m z_j,\qquad H_{3m+j}=\rho^m M_j.
\]
Every trial at the corresponding phase is exactly $\rho^m$ times its listed base trial. All comparison quantities scale by $\rho^m$, preserving every decision, not only acceptance at the selected length. Thus $x_k\to0$ and $G(x_k)\downarrow0$.

This limit is nonstationary: $d_* =(-1,16)^T$ satisfies
\[
 g_0^Td_*=-1,\quad g_1^Td_*=-1,\quad g_2^Td_*=-c+16d<-1.
\]
In particular $G(td_*)=-t$ for $t\ge0$ and $0\notin\operatorname{conv}\{g_0,g_1,g_2\}$.

% SOURCE_END D.span.5


% SOURCE_BEGIN D.span.6
\subsection{A simple coercive global failure}
The same trajectory is a deterministic counterexample for the finite polyhedral function
\[
 \overline G(u,v)=\max\{u,-u-v/8,cu+dv,v-17\}.
\]
All trial values of $G$ are positive while $v-17<0$, by~\eqref{D:trialbounds} and scaling. The fourth piece never affects the line search. The convex combination of branches $0,1,3$ with weights $(8/17,8/17,1/17)$ is identically $-1$. Equality fixes the unique minimizer $(-1,16)$, and branch $2$ is strictly below $-1$ there. Sublevels are bounded, because $u\le t$, $v\le t+17$, and $-u-v/8\le t$ give lower bounds on both coordinates. This supplies a coercive deterministic failure, but does not alone refute local convergence at its true minimizer. The next construction does.


\setcounter{section}{3}
\setcounter{theorem}{1}
\setcounter{lemma}{1}
\setcounter{corollary}{0}
\setcounter{proposition}{0}
\setcounter{remark}{0}
\setcounter{equation}{7}
% SOURCE_END D.span.6


% SOURCE_BEGIN D.span.7
\section{Two smooth convex scalar functions}
Set
\[
 \vartheta=\frac1{16},\qquad a_n=\vartheta^n\ (n\ge0),\qquad
 \epsilon_0=\frac1{128},\qquad b_* =2c-32d\in(33,34).
\]
We define scalar functions only for $b\in\{2,b_*\}$. Let
\[
 c_b=\frac{1-(1+\vartheta+\vartheta^2)/b}{1-\vartheta^2}.
\]
Define a monotone $C^1$ ramp $P_b$ on $[0,1]$ by
\[
P_b(\xi)=
\begin{cases}
0,&\xi\le c_b-\epsilon_0,\\
3t^2-2t^3,&c_b-\epsilon_0<\xi<c_b+\epsilon_0,
\quad t=(\xi-c_b+\epsilon_0)/(2\epsilon_0),\\
1,&\xi\ge c_b+\epsilon_0.
\end{cases}
\]
Then $0\le P_b'\le96$ and $\int_0^1P_b(\xi)\,d\xi=1-c_b$. Its primitive $J_b(\xi)=\int_0^\xi P_b(t)\,dt$ is
\[
J_b(\xi)=\begin{cases}
0,&\xi\le c_b-\epsilon_0,\\
2\epsilon_0(t^3-t^4/2),&c_b-\epsilon_0<\xi<c_b+\epsilon_0,\\
\xi-c_b,&\xi\ge c_b+\epsilon_0.
\end{cases}
\]
For $x\in[b\vartheta a_n,ba_n]$, put
\[
 \xi=\frac{x/(ba_n)-\vartheta}{1-\vartheta},
\]
and define
\begin{equation}\label{D:scalar}
 \varphi_b(x)=(\vartheta a_n)^3+
 ba_n^3(1-\vartheta)
 \left[\vartheta^2\xi+(1-\vartheta^2)J_b(\xi)\right].
\end{equation}
Complete its definition by
\[
\varphi_b(x)=\begin{cases}
1-3(x+1)+3(x+1)^2,&x<-1,\\
-x^3,&-1\le x\le0,\\
1+x-b,&x\ge b.
\end{cases}
\]
These intervals cover the real line, with agreeing formulas at shared endpoints.

\begin{lemma}\label{D:scalarlemma}
Each $\varphi_b$ is globally convex, $C^2$, nonnegative, coercive, and uniquely minimized at $0$. Moreover,
\begin{gather}
 \varphi_b(ba_n)=a_n^3,\quad
 \varphi_b'(ba_n)=a_n^2,\quad\varphi_b''(ba_n)=0,\label{D:jets}\\
 0\le\varphi_b''(x)\le\max\{6,102/b\}\quad(x\in\R),\label{D:curvature}\\
 \boxed{\varphi_b(ba_n+h)=a_n^3+a_n^2h\quad\text{for }|h|\le a_n/2.}\label{D:affineplateau}
\end{gather}
\end{lemma}
\begin{proof}
Differentiate~\eqref{D:scalar} on its positive intervals:
\[
 \varphi_b'(x)=a_n^2[\vartheta^2+(1-\vartheta^2)P_b(\xi)],
 \qquad
 \varphi_b''(x)=\frac{a_n(1+\vartheta)}bP_b'(\xi).
\]
The latter is nonnegative and at most $102a_n/b$. The identity
\[
 \vartheta^2+(1-\vartheta^2)(1-c_b)
 =\frac{1+\vartheta+\vartheta^2}{b}
\]
ensures the required integral and value matching at every node. The derivative is constant near each interval endpoint; thus the first two derivatives match adjoining intervals.

For $b\vartheta a_n\le x\le ba_n$, $a_n\le x/(b\vartheta)$. Consequently as $x\downarrow0$,
\[
 \varphi_b(x)=O(x^3),\qquad\varphi_b'(x)=O(x^2),\qquad
 \varphi_b''(x)=O(x).
\]
They match the negative formula at zero, including two derivatives. The joins at $-1$ and $b$ are also $C^2$. The derivative is negative on $x<0$ and positive on $x>0$; both outer extensions diverge to infinity. This proves convexity, the minimum assertions, coercivity, and~\eqref{D:curvature}.

The constant-derivative lengths below and above $ba_n$, divided by $a_n$, are respectively
\[
 \ell_b=b(1-\vartheta)(1-c_b-\epsilon_0),\qquad
 u_b=\frac{b(1-\vartheta)(c_b-\epsilon_0)}{\vartheta}.
\]
At $n=0$ the affine outer extension supplies the upper length. Exact rational interval bounds give
\[
\begin{array}{c|ccc}
 b&c_b&\ell_b&u_b\\\hline
2&0.468627450980&0.981675091912&13.824448529412\\
b_*&0.971910573260&0.635765118557&483.655089309488
\end{array}
\]
Both transition endpoints are strictly inside $(0,1)$ and both lengths exceed $1/2$. This proves~\eqref{D:affineplateau}.
\end{proof}

\paragraph{Additional regularity.}
The construction even gives globally Lipschitz second derivatives. Within a ramp,
\[
 |\varphi_b'''(x)|
 \le\frac{1+\vartheta}{b^2(1-\vartheta)}
       \frac{3}{2\epsilon_0^2}
 =\frac{139264}{5b^2}.
\]
The second derivative is continuous across every join and at the accumulation point $0$. The uniform bound above, together with the bound $6$ on the negative cubic part, proves global Lipschitz continuity of $\varphi_b''$ by summing over interval lengths. Thus the branches below are in $C^{2,1}$ as well. The first three branches are not analytic, and no branch is uniformly strongly convex.


\setcounter{section}{4}
\setcounter{theorem}{1}
\setcounter{lemma}{2}
\setcounter{corollary}{0}
\setcounter{proposition}{0}
\setcounter{remark}{0}
\setcounter{equation}{11}
% SOURCE_END D.span.7


% SOURCE_BEGIN D.span.8
\section{The single objective with bad starts arbitrarily near its minimizer}
Define
\begin{equation}\label{D:fullobjective}
\begin{split}
 f_0(u,v)&=\varphi_2(u),\\
 f_1(u,v)&=\varphi_2(-u-v/8),\\
 f_2(u,v)&=\varphi_{b_*}(cu+dv),\\
 f_3(u,v)&=v,\qquad F(u,v)=\max_{0\le i\le3}f_i(u,v).
\end{split}
\end{equation}
Each branch is convex and $C^2$. For $i\le2$, its Hessian is the scalar second derivative times $g_i g_i^T$. Thus
\[
 \|\nabla^2f_0\|\le51,\qquad
 \|\nabla^2f_1\|\le51\frac{65}{64}<52,
\]
while $b_*>33$ gives $\varphi_{b_*}''\le6$ and $c^2+d^2<2$, so
\[
 \|\nabla^2f_2\|<12,\qquad\nabla^2f_3=0.
\]
A common global upper bound $64I$ therefore holds. The preceding derivative estimate and $\|g_i\|<2$ also give a common finite Lipschitz bound for the branch Hessians (for example $60000$ suffices).

The first three branches are nonnegative and vanish at zero. If $F(x)=0$, its first two scalar arguments must vanish; $g_0$ and $g_1$ are independent, so $x=0$. The same two scalar arguments lie in bounded scalar sublevels whenever $F(x)$ is bounded above, proving compactness of every sublevel of $F$.

At zero, the first three branches have zero gradient. Therefore
\[
 F'(0;h)=\max\{0,h_2\},
\]
which proves that the unique minimizer is nonsmooth.

Put
\[
 p_n=a_n(2,-32)^T.
\]
Then $g_0^Tp_n=g_1^Tp_n=2a_n$ and $g_2^Tp_n=b_*a_n$. Lemma~\ref{D:scalarlemma} yields the exact local representation
\begin{equation}\label{D:localcone}
 \boxed{F(p_n+h)=a_n^3+a_n^2G(h)}
\end{equation}
whenever $|g_i^Th|\le a_n/2$ for $i=0,1,2$. Indeed, all three transformed branches are in their affine intervals and have value at least $a_n^3/2$. The first two slab inequalities imply $|h_1|\le a_n/2$ and $|h_2|\le8a_n$, so $f_3(p_n+h)<0$ and is inactive.

At $p_n$, the direction $d_*=(-1,16)^T$ strictly decreases all three active branches. Hence
\[
 F(p_n)=a_n^3>0,\qquad
 F'(p_n;d_*)=-a_n^2<0,
 \qquad0\notin\partial F(p_n).
\]


\setcounter{section}{5}
\setcounter{theorem}{1}
\setcounter{lemma}{2}
\setcounter{corollary}{0}
\setcounter{proposition}{0}
\setcounter{remark}{0}
\setcounter{equation}{13}
% SOURCE_END D.span.8


% SOURCE_BEGIN D.span.9
\section{All rejected trials also stay in the local model}
For each $n\ge1$, initialize
\[
 \boxed{x_0^{(n)}=p_n+a_n^2z_0,\qquad H_0=I.}
\]
The claimed deterministic trajectory is
\begin{equation}\label{D:finalorbit}
 x_{3m+j}^{(n)}=p_n+a_n^2\rho^mz_j,\qquad
 H_{3m+j}^{(n)}=\rho^mM_j,\qquad j=0,1,2.
\end{equation}
At the corresponding phase, \emph{every} trial generated by the line search has the form
\[
 p_n+a_n^2\rho^m z_{\rm trial},
\]
where $z_{\rm trial}$ is one of the nine base trials already certified. By~\eqref{D:trialbounds},
\[
 |g_i^T(a_n^2\rho^mz_{\rm trial})|
 <3a_n^2\le\frac{3a_n}{16}<\frac{a_n}{2}.
\]
Thus~\eqref{D:localcone} holds at every accepted and rejected trial. Their values and gradients are
\[
 F(p_n+a_n^2\rho^mz_{\rm trial})
 =a_n^3+a_n^4\rho^mG(z_{\rm trial}),
 \qquad \nabla F=a_n^2g_{\rm active}.
\]
Every function difference and every predicted directional decrease is the corresponding seed quantity multiplied by $a_n^4\rho^m$. The additive $a_n^3$ cancels from Armijo. Therefore all rejection and acceptance decisions are the same as in the seed table. In particular the accepted lengths are outputs of the deterministic search.

Substituting gradients scaled by $a_n^2$ and steps scaled by $a_n^2\rho^m$ into~\eqref{D:bfgs} leaves the projection factors unchanged and scales the rank-one term by $\rho^m$. This proves the matrix part of~\eqref{D:finalorbit}. The point and matrix cycle closures complete the induction. All finite matrices are SPD, and all oracle evaluations are at differentiable points.

Consequently,
\[
 x_k^{(n)}\to p_n\ne0,\qquad
 F(x_k^{(n)})\downarrow a_n^3>0.
\]
Yet $x_0^{(n)}\to0$ and $F(x_0^{(n)})\to0$. Also the supremum of the norms of \emph{all accepted and trial points} in the $n$th trajectory is bounded by
\[
 a_n\|(2,-32)\|+a_n^2\max_{\rm base\ trials}\|z_{\rm trial}\|\to0.
\]
This proves Theorem~\ref{D:main}. The failure occurs entirely within the local model.


\endgroup
% SOURCE_END D.span.9


\section*{Relation to earlier work}
The following comparisons are bibliographic scholarship, separate from the formal theorem assertions.

The unit-first search in this chapter is Lewis--Overton's Algorithm 4.6 \cite{LO}. Earlier negative results include their convex subgradient/exact-search construction with null steps \cite{NEG_LO08}, smooth globally nonconvex examples \cite{NEG_M04,NEG_M07,NEG_M14,NEG_DAI02,NEG_DAI13}, Clarke-stationary limits with descent directions \cite{NEG_LZ15}, and scaled memoryless BFGS failure \cite{NEG_AO21}. Those settings differ from the convex, differentiable-query, full-update statement here. The recent planar counterexample of Diao uses a smooth nonconvex objective and an escaping exact-search trajectory \cite{NEG_DIAO26}. No exact match for the complete theorem was identified in the documented search.

Prescribing an algebraic orbit and then constructing an objective that realizes its update equations has substantial precedents in Mascarenhas's divergence constructions \cite{NEG_M04,NEG_M07,NEG_M14}. The convex interpolation and all-query containment in this chapter are specific additional properties. Preconditioned residual maps and rigorous root enclosures are standard validated-numerics tools; compare Rump \cite[Theorem 13.3]{NEG_RUMP10}. Here the required rational inequalities have their own Lean proofs.
